%! Describing the Distribution of a Quantitative Variable — Unit 1, Lesson 1.4 — Warm-Up (Answer Key)
% ONE PAGE, blank and key. Three items, set at 12pt like every other student component.
% Do not add a fourth item — the page is full.
% GRADUAL RELEASE: the warm-up activates prior knowledge before the guided notes. Item 1 is
% the spiral back to Lesson 1.3 (two interval widths, one data set); item 2 is the
% count-to-the-middle move notes section 1 needs, so finding a center is not the obstacle
% there; item 3 puts a lone far-off value in front of them in plain language — notes section 2
% names it an OUTLIER, and the empty stretch before it a GAP, fifteen minutes later.
% NO SKETCH QUESTIONS: the item 3 display is pre-drawn; students only read it.
% Mirrors the blank exactly apart from the package swap, the pageheader, and the answers.
% NO teachernote here — teacher prose lives in the lesson plan.
\documentclass[12pt]{article}
\usepackage{apstats-article}
\usepackage{apstats-key}   % pulls in -boxes; do NOT also load -boxes

\begin{document}

\pageheader{Unit 1, Lesson 1.4}{Warm-Up --- Answer Key}

\vspace{0.20in}

\begin{enumerate}[leftmargin=1.9em, itemsep=0.42in]

  \item \textbf{From last lesson.} A biologist measured the lengths of 40 fish and drew the
        data twice. With blocks 2 cm wide the histogram showed two separate peaks; with blocks
        10 cm wide the very same 40 lengths showed one single hump.

        \vspace{10pt}
        Which of the two pictures is \emph{wrong}? \ans{Neither --- both are correct}

        \vspace{14pt}
        What should you check before you believe what any histogram looks like?
        \ans{the width of the blocks}

  \item Nine pumpkins were weighed at a county fair. Their weights, in pounds, are already
        written in order from lightest to heaviest.

        \vspace{10pt}
        \begin{center}
        4 \quad 6 \quad 6 \quad 7 \quad 9 \quad 9 \quad 11 \quad 14 \quad 16
        \end{center}
        \vspace{10pt}

        How many weights are listed? \ans{9} \qquad
        Counting in from both ends, which value lands exactly in the middle? \ans{9 lb}

  \item Fifteen students reported how many books they read over the summer.

        \vspace{4pt}
        \begin{center}
        \begin{tikzpicture}[scale=0.9]
          \draw[->, thick] (-0.2,0) -- (8.6,0);
          \foreach \x/\lab in {0/0, 1/2, 2/4, 3/6, 4/8, 5/10, 6/12, 7/14, 8/16} {
            \draw (\x,-0.07) -- (\x,0.07); \node[below, font=\tiny] at (\x,-0.1) {\lab};}
          \foreach \x/\n in {1/2, 1.5/3, 2/4, 2.5/2, 3/2, 3.5/1, 7.5/1} {
            \foreach \k in {1,...,\n} { \fill[navy] (\x,{0.22*\k}) circle (0.075); }}
          \node[font=\tiny] at (4.2,-0.85) {Books read over the summer};
        \end{tikzpicture}
        \end{center}
        \vspace{4pt}

        Most of the students read between \ans{2} and \ans{7} books. \quad
        One student sits far away from all the others, at \ans{15} books.

        \vspace{10pt}
        In one sentence, describe where the other fourteen students are compared with that one.\\[10pt]
        \ansline{The other fourteen are bunched down at the low end, between 2 and 7 books,}\\[14pt]\ansline{and then nothing at all until that single student out at 15 books.}

\end{enumerate}

\end{document}
